\subsubsection{Subproceso 1: Acreditación y Pre-requisitos}

El propósito de este subproceso es verificar que el estudiante regular cumpla con las condiciones administrativas, de pagos y académicas antes de iniciar su inscripción de cursos. En esta etapa se comprueba el pago del autoseguro médico, la ausencia de deudas en la Facultad y el cierre de notas del ciclo anterior para calcular su promedio ponderado, definir su límite de créditos y asignarle un turno de matrícula.

Este subproceso abre la matrícula regular y funciona como filtro inicial. Si el alumno cumple con todos los requisitos, el sistema actualiza su estado a «Habilitado» y genera su ficha de matrícula con los cursos que puede llevar según su plan de estudios, permitiendo así que pase al Subproceso 2 en la fecha y hora de su turno.
